% === ЗАДАЧА 7 ===
% Тег: #ОПТИКА #ВОЛНОВАЯ_ОПТИКА #ДИФРАКЦИОННАЯ_РЕШЕТКА #КИМ_23 #ВАРИАНТ_09
\begin{taskbasic}[code={\label{task:wave_opt_v9_n23}}]
На дифракционную решетку, имеющую $500\text{ штрихов}$ на $1\text{ мм}$, падает по нормали параллельный пучок монохроматического света с длиной волны $\lambda$. Максимум второго порядка наблюдается под углом $\varphi = 36{,}9^\circ$ ($\sin 36{,}9^\circ = 0{,}6$). Найдите длину волны $\lambda$. Ответ выразите в нанометрах (нм).

\kim{23}
\end{taskbasic}
\answer{600}

% === ЗАДАЧА 8 ===
% Тег: #ОПТИКА #ВОЛНОВАЯ_ОПТИКА #ДИФРАКЦИОННАЯ_РЕШЕТКА #КИМ_23 #ВАРИАНТ_10
\begin{taskbasic}[code={\label{task:wave_opt_v10_n23}}]
На дифракционную решетку падает по нормали параллельный пучок монохроматического света с длиной волны $\lambda = 500\text{ нм}$. Найти период дифракционной решетки, если дифракционный максимум третьего порядка наблюдается под углом, синус которого равен $0{,}6$. Ответ выразите в микрометрах (мкм).

\kim{23}
\end{taskbasic}
\answer{2,5}

% === ЗАДАЧА 1 ===
% Тег: #ОПТИКА #ВОЛНОВАЯ_ОПТИКА #ДИФРАКЦИОННАЯ_РЕШЕТКА #КИМ_23
\begin{taskbasic}[code={\label{task:wave_opt_diffraction_grating_max_k_300}}]
На дифракционную решётку, имеющую $300\text{ штрихов}$ на $1\text{ мм}$, перпендикулярно её поверхности падает узкий луч монохроматического света частотой $5{,}6 \cdot 10^{14}\text{ Гц}$. Каков максимальный порядок дифракционного максимума, доступного для наблюдения?

\kim{23}
\end{taskbasic}
\answer{6}

% === ЗАДАЧА 2 ===
% Тег: #ОПТИКА #ВОЛНОВАЯ_ОПТИКА #ДИФРАКЦИОННАЯ_РЕШЕТКА #КИМ_23
\begin{taskbasic}[code={\label{task:wave_opt_diffraction_grating_max_k_500}}]
На дифракционную решётку, имеющую $500\text{ штрихов}$ на $1\text{ мм}$, перпендикулярно её поверхности падает узкий луч монохроматического света частотой $5 \cdot 10^{14}\text{ Гц}$. Каков максимальный порядок дифракционного максимума, доступного для наблюдения?

\kim{23}
\end{taskbasic}
\answer{3}